\documentclass{article}
\usepackage[T1]{fontenc}
\usepackage{lmodern}
\usepackage{graphicx}
\usepackage{amsmath,amssymb}
\usepackage[a4paper,margin=2cm]{geometry}
\usepackage[absolute,overlay]{textpos}
\setlength{\TPHorizModule}{1mm}
\setlength{\TPVertModule}{1mm}
\usepackage[indonesian]{babel}
\usepackage{hyperref}
\setlength{\emergencystretch}{2em}
% unit: Halaman 4

\begin{document}

% === Halaman 4 (python/native) ===

• Keamanan algoritma ini terletak pada sulitnya menghitung logaritma 

diskrit.

• Persoalan logaritma diskrit: Jika p adalah bilangan prima dan g dan y 

adalah sembarang bilangan bulat, carilah x sedemikian sehingga 



gx  y (mod p)

 Contoh:   7x  15 (mod 41), berapakah nilai x?

• Sebelum membahas algoritma ElGamal lebih lanjut, maka kita perlu 

memahami terlebih dahulu  tentang logaritma diskrit dan akar primitif 

dari sebuah bilangan prima.

4

\end{document}
